\begin{figure}[!htbp]
\centering
\begin{adjustbox}{max width=\linewidth}
\begin{tikzpicture}[x=0.72cm, y=0.72cm]
  \foreach \x/\y in {20/10, 21/11, 22/12, 23/12, 24/13, 25/14, 26/15, 27/16, 28/16, 29/17, 30/18}
    \filldraw[fill=fwine!45, draw=fgink!60, line width=0.3pt] ({\x-20-0.5},{\y-10-0.5}) rectangle ++(1,1);
  \draw[fgink!35, line width=0.3pt] (-0.5,-0.5) grid[step=1] (10.5,8.5);
  \draw[fgink, line width=0.7pt] (0,0) -- (10,8);
  \foreach \x in {20,22,...,30} \node[font=\scriptsize, below] at ({\x-20},-0.5) {\x};
  \foreach \y in {10,12,...,18} \node[font=\scriptsize, left] at (-0.5,{\y-10}) {\y};
  \node[lbl, anchor=west, text width=40mm] at (11.0,4.0) {shaded pixels chosen by Bresenham's algorithm; the true line is drawn over them};
\end{tikzpicture}
\end{adjustbox}
\caption{Bresenham's algorithm for the line from (20, 10) to (30, 18): at each column it picks the pixel
nearer the true line using only integer arithmetic.}
\end{figure}
\begin{center}
\renewcommand{\arraystretch}{1.2}
\begin{tabular}{c r c}
\toprule
$k$ & $p_k$ & next pixel \\
\midrule
0 & 6 & (21, 11) \\
1 & 2 & (22, 12) \\
2 & $-2$ & (23, 12) \\
3 & 14 & (24, 13) \\
4 & 10 & (25, 14) \\
5 & 6 & (26, 15) \\
\bottomrule
\end{tabular}
\end{center}
